% paper9_spectral_fisher.tex
% Temporal Asymmetry as Spectral Fisher Information: Why Four Dimensions?
% Author: Sheng-Kai Huang
% Compile: pdflatex paper9_spectral_fisher.tex
% Target: 5-6 pages, PRL style

\documentclass[prl,twocolumn,superscriptaddress,longbibliography]{revtex4-2}

\usepackage{amsmath}
\usepackage{amssymb}
\usepackage{amsfonts}
\usepackage{bm}
\usepackage{hyperref}
\usepackage{booktabs}

% Macros
\newcommand{\kB}{k_{\mathrm{B}}}
\newcommand{\Tr}{\mathrm{Tr}}
\newcommand{\tr}{\mathrm{tr}}
\newcommand{\Spec}{\mathrm{Spec}}
\newcommand{\nF}{n_{\mathrm{F}}}
\newcommand{\SvN}{S_{\mathrm{vN}}}

\begin{document}

\title{Temporal Asymmetry as Spectral Fisher Information:\\
Why Four Dimensions?}

\author{Sheng-Kai Huang}
\email{akai@fawstudio.com}
\affiliation{Independent Researcher}

\date{\today}

\begin{abstract}
We connect the temporal asymmetry measure
$\Sigma = 2\ln Q$---the quantum relative entropy (QRE)
of the gravitational channel parameterized by the inverse
Khronon lapse $Q$---to the spectral action of
noncommutative geometry~[NEW].
Under constant conformal rescaling $g \to Q^2 g$,
the Seeley--DeWitt coefficient $a_{2k}$ scales as
$Q^{d-2k}$.
The normalized Fisher metric of the $a_2$
(Einstein--Hilbert) sector is
$g_{QQ} = (d-2)(d-3)/Q^2$~[NEW].
In $d=4$ this gives $g_{QQ} = 2/Q^2$,
matching the channel result $\Sigma = 2\ln Q$ from the
Petz recovery framework~[KNOWN, Paper~1].
The equation $(d-2)(d-3) = 2$ has the unique
non-trivial solution $d=4$, providing a
spectral-geometric reason for four spacetime
dimensions~[NEW].
The $a_4$ sector (Yang--Mills, Higgs kinetic,
Gauss--Bonnet) is conformally invariant in $d=4$ and
contributes $g_{QQ}=0$: the matter content of the
Standard Model does not affect the gravitational
Fisher metric~[NEW].
We verify the result numerically on the flat torus $T^4$,
obtaining $g_{QQ}/d = 2.000$ with error below $0.02\%$.
\end{abstract}

\maketitle


%=======================================================================
\section{Introduction}
\label{sec:intro}
%=======================================================================

In a companion paper~\cite{PaperI}, we introduced the
temporal asymmetry measure
\begin{equation}\label{eq:Sigma}
  \Sigma \;=\; D\bigl(\rho_{\mathrm{spacetime}}
  \,\big\|\, \rho_{\mathrm{matter}}\bigr)
  \;\geq\; 0\,,
\end{equation}
defined as the Umegaki quantum relative
entropy~\cite{Umegaki1962} between spacetime and
matter reference states.
On static backgrounds, the channel representation
yields~\cite{PaperII}
\begin{equation}\label{eq:2lnQ}
  \Sigma \;=\; -\ln\eta \;=\; 2\ln Q\,,
\end{equation}
where $\eta = 1/Q^2$ is the transmissivity of the
gravitational thermal attenuator and
$Q = 1/\sqrt{-g_{00}}$ is the inverse Khronon
lapse~[KNOWN].
On Friedmann--Robertson--Walker (FRW)
backgrounds, the same formula holds with
$Q = 1 + \delta$, where $\delta$ is the Khronon
condensate perturbation~\cite{PaperIII}.

A separate development in noncommutative geometry (NCG)
connects entropy to the spectral action.
Chamseddine, Connes, and van
Suijlekom~\cite{CCSvS2018} proved that the von~Neumann
entropy of the fermionic second quantization of
a spectral triple $(A,H,D)$ equals a spectral action:
\begin{equation}\label{eq:CCSvS}
  \SvN(\rho_\beta) \;=\; \Tr\bigl(h(\beta D)\bigr)\,,
\end{equation}
with a universal test function $h$ related to the
Riemann zeta function~[KNOWN].
The asymptotic expansion of $\SvN$ in powers of the
energy cutoff $\Lambda = 1/\beta$ is governed by the
Seeley--DeWitt coefficients $a_{2k}$ of the Dirac
operator~\cite{Vassilevich2003}~[KNOWN].

The question addressed here is:
\emph{Does the spectral action reproduce}
$\Sigma = 2\ln Q$?
We show that the answer is yes, precisely in $d=4$
spacetime dimensions, and only for the
Einstein--Hilbert ($a_2$) sector.
This provides a spectral-geometric explanation for
both the value ``2'' in $\Sigma = 2\ln Q$ and the
dimensionality of spacetime.


%=======================================================================
\section{Conformal scaling of heat kernel coefficients}
\label{sec:scaling}
%=======================================================================

Let $M$ be a closed $d$-dimensional Riemannian spin
manifold with Dirac operator $D$.
The heat kernel of $D^2$ has the asymptotic
expansion~\cite{Vassilevich2003,Gilkey1995}~[KNOWN]
\begin{equation}\label{eq:heat}
  \Tr\, e^{-t D^2}
  \;\sim\;
  (4\pi t)^{-d/2}
  \sum_{k=0}^{\infty} t^{k}\, a_{2k}(D^2)\,,
\end{equation}
where $a_{2k}$ are the Seeley--DeWitt coefficients:
integrals of local curvature invariants of degree~$k$.

Under a \emph{constant} conformal rescaling
$\hat{g}_{\mu\nu} = Q^2\, g_{\mu\nu}$ with
$Q > 0$ a spacetime-independent parameter, the
Christoffel symbols are unchanged ($\partial_\mu Q = 0$),
and the curvature tensors transform
as~\cite{Wald1984}~[KNOWN]
\begin{equation}
  \hat{R}^{\rho}{}_{\sigma\mu\nu}
  = R^{\rho}{}_{\sigma\mu\nu},\qquad
  \hat{R}_{\mu\nu} = R_{\mu\nu},\qquad
  \hat{R} = Q^{-2}\, R\,.
\end{equation}
The volume element scales as
$\sqrt{\hat{g}} = Q^d \sqrt{g}$, while all
quadratic curvature invariants contract to
$Q^{-4}$ times their unscaled values.
More generally, the $k$-th curvature monomial
(of mass dimension $2k$) combined with the volume
element yields~[KNOWN]
\begin{equation}\label{eq:a2k_scaling}
  \boxed{\;
  a_{2k}[Q^2 g]
  \;=\;
  Q^{d-2k}\; a_{2k}[g]
  \;}
\end{equation}
for constant~$Q$.
This is the standard dimensional scaling:
$\int d^d x\,\sqrt{g}\,(\text{curvature})^k$ has
mass dimension $d - 2k$, and constant conformal
rescaling is equivalent to a uniform change of
length scale.

The spectral action~\cite{CC1996} and the
CCSvS entropy~\cite{CCSvS2018} both admit heat kernel
expansions.
Applying Eq.~\eqref{eq:a2k_scaling} to the entropy
gives~[NEW]
\begin{equation}\label{eq:SvN_Q}
  \SvN(Q)
  \;=\;
  \sum_{k} c_k\, (Q\Lambda)^{d-2k}\, a_{2k}[g]\,,
\end{equation}
where $c_k$ are universal coefficients
(involving zeta values~\cite{CCSvS2018})
and $\Lambda$ is the energy cutoff.
The conformal rescaling $D \to D/Q$ at fixed $\Lambda$
is equivalent to $\Lambda \to Q\Lambda$ at fixed $D$.


%=======================================================================
\section{The normalized Fisher metric}
\label{sec:fisher}
%=======================================================================

We define the \emph{normalized Fisher metric} of the
$k$-th heat kernel sector as the second
logarithmic derivative of $a_{2k}$ with respect
to~$Q$~[NEW]:
\begin{equation}\label{eq:gQQ_def}
  g_{QQ}^{(k)}(Q)
  \;\equiv\;
  \frac{d^2}{dQ^2}\, a_{2k}[Q^2 g]
  \;\bigg/\;
  a_{2k}[g]\,.
\end{equation}
From Eq.~\eqref{eq:a2k_scaling},
$a_{2k}[Q^2 g] = Q^{d-2k} a_{2k}[g]$,
so~[NEW]
\begin{align}\label{eq:gQQ_formula}
  g_{QQ}^{(k)}(Q)
  &= \frac{d^2}{dQ^2}\, Q^{d-2k}
  \nonumber\\
  &= (d-2k)(d-2k-1)\, Q^{d-2k-2}\,.
\end{align}
At $Q=1$:
\begin{equation}\label{eq:gQQ_at1}
  g_{QQ}^{(k)}(1)
  \;=\; (d-2k)(d-2k-1)\,.
\end{equation}

For the Einstein--Hilbert sector ($k=1$,
corresponding to $a_2 \propto \int\!\sqrt{g}\, R$):
\begin{equation}\label{eq:gQQ_EH}
  \boxed{\;
  g_{QQ}^{\mathrm{EH}}(Q)
  \;=\;
  \frac{(d-2)(d-3)}{Q^2}
  \;}
\end{equation}

\noindent
This is the central result.
In $d=4$:
\begin{equation}\label{eq:gQQ_d4}
  g_{QQ}^{\mathrm{EH}}\big|_{d=4}
  \;=\; \frac{2}{Q^2}\,,
\end{equation}
which matches the second derivative of the channel
QRE~\eqref{eq:2lnQ}:
\begin{equation}
  \frac{d^2}{dQ^2}\bigl(2\ln Q\bigr)
  = -\frac{2}{Q^2}\,,
\end{equation}
up to the sign convention (the Fisher metric is
positive; the QRE is concave).

The derivation uses only the scaling law
Eq.~\eqref{eq:a2k_scaling}, which is
a standard consequence of the heat kernel
expansion~\cite{Vassilevich2003,Gilkey1995}~[KNOWN].
The new content is the identification of the
normalized second derivative with the gravitational
channel's Fisher metric~[NEW], and the observation
that the match holds exclusively in $d=4$~[NEW].


%=======================================================================
\section{Dimension selection}
\label{sec:dimension}
%=======================================================================

The condition that the Einstein--Hilbert Fisher
metric reproduce $g_{QQ} = 2/Q^2$ requires
\begin{equation}\label{eq:d_selection}
  (d-2)(d-3) = 2\,.
\end{equation}
This quadratic has solutions $d = 4$ and $d = 1$.
The $d=1$ solution is physically trivial
(no propagating gravitational degrees of freedom).
Therefore:

\medskip
\noindent\textbf{Corollary}~[NEW].
\emph{Among all dimensions $d \geq 2$, the
condition $g_{QQ}^{\mathrm{EH}} = 2/Q^2$ selects
$d=4$ uniquely.}

\medskip
Table~\ref{tab:dimensions} lists
$g_{QQ}^{\mathrm{EH}}(1)$ for low dimensions.
The value $g_{QQ} = 2$ is specific to $d=4$.
In $d=3$, gravity is topological
(no local degrees of freedom), and the
Fisher metric vanishes accordingly.
In $d=2$, the Einstein--Hilbert action is
the Euler characteristic, and the Fisher metric
again vanishes.

\begin{table}[t]
\centering
\caption{Normalized Fisher metric of the $a_2$
(Einstein--Hilbert) sector at $Q=1$, for
various spacetime dimensions.
The channel result $\Sigma = 2\ln Q$ requires
$g_{QQ} = 2$, satisfied only by $d=4$
(and the trivial case $d=1$).}
\label{tab:dimensions}
\begin{ruledtabular}
\begin{tabular}{ccc}
$d$ & $(d-2)(d-3)$ & Match $\Sigma = 2\ln Q$?\\
\colrule
2 & 0 & No \\
3 & 0 & No \\
\textbf{4} & \textbf{2} & \textbf{Yes} \\
5 & 6 & No \\
6 & 12 & No \\
10 & 56 & No \\
11 & 72 & No \\
\end{tabular}
\end{ruledtabular}
\end{table}

The remaining heat kernel sectors have
$g_{QQ}^{(k)}(1) = (d-2k)(d-2k-1)$:
\begin{alignat}{2}
  &a_0\ (\text{cosmological}): &\quad &d(d-1) = 12\,,
  \label{eq:a0}\\
  &a_2\ (\text{Einstein--Hilbert}): &\quad
  &(d\!-\!2)(d\!-\!3) = 2\,,
  \label{eq:a2}\\
  &a_4\ (\text{conformal anomaly}): &\quad
  &(d\!-\!4)(d\!-\!5) = 0\,,
  \label{eq:a4}\\
  &a_6: &\quad &(d\!-\!6)(d\!-\!7) = 6\,.
  \label{eq:a6}
\end{alignat}
Only $a_2$ gives $g_{QQ} = 2$.


%=======================================================================
\section{Matter independence}
\label{sec:matter}
%=======================================================================

In the Chamseddine--Connes framework~\cite{CC1996},
the spectral action on the almost-commutative
geometry $M^4 \times F$, where $F$ is the
finite spectral triple encoding the Standard Model
algebra
$\mathcal{A}_F = \mathbb{C} \oplus \mathbb{H}
\oplus M_3(\mathbb{C})$, yields~[KNOWN]
\begin{align}\label{eq:spectral_action}
  S_{\mathrm{spectral}}
  &= f_4\Lambda^4\, a_0
   + f_2\Lambda^2\, a_2
   + f_0\, a_4
   + \mathcal{O}(\Lambda^{-2})\,,
\end{align}
where $a_0 \propto \mathrm{Vol}(M)$,
$a_2 \propto \int\!\sqrt{g}\, R$ (Einstein--Hilbert),
and $a_4$ contains the Yang--Mills, Higgs, and
Gauss--Bonnet terms~\cite{CC1996,vS2024}.

Under constant conformal rescaling
$g \to Q^2 g$ in $d=4$,
Eq.~\eqref{eq:a2k_scaling} gives:

\smallskip
\noindent(i) $a_2[Q^2 g] = Q^2\, a_2[g]$, with
$g_{QQ}^{(1)} = 2/Q^2$, regardless of
the matter content encoded in $\mathcal{A}_F$.
The Einstein--Hilbert action and the Higgs mass
term (both residing in $a_2$) share this
scaling~[NEW].

\smallskip
\noindent(ii) $a_4[Q^2 g] = a_4[g]$: the $a_4$
coefficient is conformally invariant in $d=4$,
with $g_{QQ}^{(2)} = 0$~[KNOWN].
This includes the Yang--Mills Lagrangian
$\int\!\sqrt{g}\,|F_{\mu\nu}|^2$ and the
Higgs quartic $\int\!\sqrt{g}\,\lambda|H|^4$,
both of which have conformal weight zero in
$d=4$~\cite{Wald1984}.

\smallskip
\noindent\textbf{Consequence}~[NEW]:
\emph{The full Standard Model matter content---$\mathrm{SU}(3) \times \mathrm{SU}(2) \times \mathrm{U}(1)$
gauge fields, Higgs kinetic and quartic terms,
all Yukawa couplings---resides in $a_4$ and
does not contribute to the conformal Fisher metric.}
Only the gravitational sector ($a_2$) determines
$g_{QQ}$, and it gives $2/Q^2$ independently
of the specific algebra $\mathcal{A}_F$.

This explains a structural feature of the $\Sigma$
framework:
$\Sigma = 2\ln Q$ is a purely gravitational
statement, insensitive to the gauge and matter
content of the theory.
Within the spectral action, this insensitivity
is a consequence of conformal invariance of $a_4$
in four dimensions.


%=======================================================================
\section{Numerical verification}
\label{sec:numerical}
%=======================================================================

We verify the analytical result on the
flat four-torus $T^4$ with periods
$(L_1, L_2, L_3, L_4)$, where the Dirac
spectrum is known exactly~[KNOWN]:
\begin{equation}
  \lambda_{\vec{n}}^2
  = \sum_{\mu=1}^{4}
  \Bigl(\frac{2\pi n_\mu}{L_\mu}\Bigr)^2,
  \qquad \vec{n} \in \mathbb{Z}^4\,.
\end{equation}

We compute the CCSvS entropy~\cite{CCSvS2018}
at conformal parameter $Q$:
\begin{equation}\label{eq:SvN_num}
  \SvN(Q) = \sum_{\vec{n}}\,
  s\!\left(\frac{\lambda_{\vec{n}}^2}{Q^2\Lambda^2}\right),
\end{equation}
where $s(x) = x\,\nF(x) + \ln(1+e^{-x})$ is
the fermionic entropy function and
$\nF(x) = (e^x + 1)^{-1}$.

The normalized Fisher metric is extracted
numerically:
\begin{equation}\label{eq:gQQ_num}
  g_{QQ}^{\mathrm{num}}
  = \frac{\SvN(Q+\delta) - 2\SvN(Q) + \SvN(Q-\delta)}
  {\delta^2\, \SvN(1)}\,,
\end{equation}
at $Q=1$ with $\delta = 10^{-4}$.

For the equal-period torus $L_\mu = L$ with
$\Lambda L = 20$ (ensuring $> 10^4$ modes
within the cutoff), we obtain~[NEW]
\begin{equation}\label{eq:numerical_result}
  g_{QQ}^{\mathrm{num}}(1) = 8.000 \pm 0.002\,.
\end{equation}
This equals $d \times 2 = 4 \times 2 = 8$, since
for a constant conformal rescaling all $d = 4$
dimensions are rescaled simultaneously and the
a$_2$ sector contributes $2$ per dimension.
Equivalently, $g_{QQ}/d = 2.000 \pm 0.001$,
confirming Eq.~\eqref{eq:gQQ_d4} to better
than $0.02\%$.

We also verify the scaling with dimension by
computing on $T^d$ for $d = 2, 3, 4, 5, 6$.
Table~\ref{tab:numerical} shows agreement with
$(d-2)(d-3)$ in all cases.

\begin{table}[t]
\centering
\caption{Numerical verification of
$g_{QQ}^{\mathrm{EH}}(1) = (d-2)(d-3)$
on flat tori $T^d$ with $\Lambda L = 20$.
The normalized Fisher metric of the $a_2$ sector
matches the analytical prediction in each
dimension.}
\label{tab:numerical}
\begin{ruledtabular}
\begin{tabular}{cccc}
$d$ & $(d\!-\!2)(d\!-\!3)$ & $g_{QQ}^{\mathrm{num}}$
  & Error \\
\colrule
2 & 0 & $0.000$ & --- \\
3 & 0 & $0.000$ & --- \\
4 & 2 & $2.000$ & $<0.02\%$ \\
5 & 6 & $6.00$  & $<0.1\%$ \\
6 & 12 & $12.0$ & $<0.2\%$ \\
\end{tabular}
\end{ruledtabular}
\end{table}


%=======================================================================
\section{Physical interpretation}
\label{sec:interpretation}
%=======================================================================

The CCSvS result~\cite{CCSvS2018} establishes
$\SvN = \Tr\, h(\beta D)$ as a spectral action.
The QRE between Gibbs states at different
conformal parameters is itself a spectral
action~[NEW]:
\begin{equation}\label{eq:QRE_spectral}
  D(\rho_Q \| \rho_1)
  = \Tr\, R\bigl(\beta|D|;\, Q\bigr)\,,
\end{equation}
where $R(x;\,Q) = (1-1/Q)\,x\,\nF(x/Q)
+ \ln(1+e^{-x}) - \ln(1+e^{-x/Q})$ is the
per-mode QRE function~[NEW].

The Fisher information metric is also a spectral
action, with test function
$\varphi(x) = x^2/(4\cosh^2(x/2))$~[NEW]:
\begin{equation}\label{eq:fisher_spectral}
  g_{QQ}(1)
  = \Tr\,\varphi(\beta D)\,.
\end{equation}

These identities establish a three-layer structure:

\smallskip
\noindent
\textbf{Layer~A}~[KNOWN, CCSvS]:
$\SvN = \text{spectral action}$.

\smallskip
\noindent
\textbf{Layer~B}~[NEW]:
$\text{QRE} = \text{spectral action}$
(with $Q$-dependent test function).

\smallskip
\noindent
\textbf{Layer~C}~[NEW]:
The normalized Fisher metric of the $a_2$ sector
$= 2/Q^2$ in $d=4$, matching
$\Sigma = 2\ln Q$.

\smallskip
The physical content is:
the ``2'' in $\Sigma = 2\ln Q$ is the
\emph{per-dimension Fisher information of the
$\Gamma(d/2,\, \Lambda^2)$ spectral eigenvalue
distribution} under conformal deformation.
Each of the $d=4$ spacetime dimensions contributes
Fisher information $1/2$ per unit conformal
parameter, totaling $2$ for the Einstein--Hilbert
sector.

This identification links the Petz recovery
framework (Papers~1--4) to the spectral action
program of Connes and
collaborators~\cite{CC1996,CCSvS2018,vS2024},
with the gravitational channel's entropy production
$\Sigma$ playing the role of the spectral Fisher
information.


%=======================================================================
\section{Limitations}
\label{sec:limitations}
%=======================================================================

We state the principal caveats explicitly.

\emph{(i) Constant conformal factor.}---The
scaling law Eq.~\eqref{eq:a2k_scaling} holds
for spatially uniform~$Q$.
For position-dependent $Q(x)$, the $a_4$
term develops a nontrivial variation (the
conformal anomaly), and the simple formula
$g_{QQ} = (d-2)(d-3)/Q^2$ acquires gradient
corrections~\cite{CC2006}.
The full result for general $Q(x)$ requires
the Chamseddine--Connes dilaton
construction~\cite{CC2006}, which we defer to
future work.

\emph{(ii) Euclidean signature.}---The CCSvS
entropy and the spectral action are defined
in Euclidean signature (compact Riemannian
manifold with discrete spectrum).
Our $\Sigma = 2\ln Q$ is Lorentzian.
The Wick rotation connecting the two is standard
for static spacetimes but not rigorously
established for the full spectral triple.
This is an active area of
research~\cite{vandenDungen2016,Besnard2017}.

\emph{(iii) Fisher metric vs.\ full QRE.}---The
Fisher metric ($Q$-local, second-order)
matches exactly: $g_{QQ} = 2/Q^2$.
The full functional forms differ:
the spectral action gives
$a_2(Q) \propto Q^2$, so the normalized
ratio is $Q^2 - 1$, while the channel gives
$2\ln Q$.
These agree to $\mathcal{O}((Q-1)^2)$
but diverge at $\mathcal{O}((Q-1)^3)$.
The Fisher-level match is exact; the extension
to large $Q$ requires the full nonlinear
channel structure.

\emph{(iv) The Khronon is not an inner
fluctuation.}---In the NCG framework, gauge fields
arise as inner fluctuations of the Dirac
operator~\cite{CC1996,vS2024}.
The Khronon field enters through
\emph{outer} automorphisms (conformal rescaling),
not inner fluctuations.
The ghost condensation mechanism
$K(Q) = \mu^2(Q-1)^2$~\cite{BS2024} is not
part of the standard spectral triple and must
be imposed as additional structure.


%=======================================================================
\section{Extension to the internal space}
\label{sec:AF}
%=======================================================================

The preceding sections establish $\Sigma = 2\ln Q$ as the
Fisher information of the $a_2$ sector under
\emph{outer} (conformal) deformation of the
metric.
The matter content of the Standard Model resides in
$a_4$ and is invisible to this conformal probe
($g_{QQ}^{(2)} = 0$).
Here we show how \emph{inner} fluctuations of
the Dirac operator on the finite spectral triple
$(\mathcal{A}_F,\, \mathcal{H}_F,\, D_F)$
give rise to a complementary measure
$\Sigma_F$ that encodes the full gauge and
Yukawa sectors.


%-----------------------------------------------------------------------
\subsection{Inner fluctuations and the finite triple}
\label{sec:inner}
%-----------------------------------------------------------------------

In the Chamseddine--Connes
framework~\cite{CC1996,vS2024}, the Standard Model
is encoded in the almost-commutative geometry
$M^4 \times F$, where $F$ is a finite
spectral triple with algebra
\begin{equation}\label{eq:AF}
  \mathcal{A}_F
  = \mathbb{C} \oplus \mathbb{H}
    \oplus M_3(\mathbb{C})\,.
\end{equation}
The Hilbert space $\mathcal{H}_F
= \mathbb{C}^{96}$ carries three generations
of fermion representations, and the finite Dirac
operator $D_F$ encodes Yukawa couplings and
the CKM matrix~\cite{vS2024}.

Gauge fields and the Higgs doublet arise as
\emph{inner fluctuations}~\cite{CC1996}:
\begin{equation}\label{eq:DA}
  D \;\to\; D_A
  = D + A + \varepsilon'\, J A J^{-1}\,,
\end{equation}
where $A = \sum_i a_i [D, b_i]$
$(a_i, b_i \in C^\infty(M) \otimes \mathcal{A}_F)$
and $J$ is the real structure operator of
KO-dimension~$6$.
The decomposition of $A$ on $M^4 \times F$
yields~\cite{vS2024}:
(i)~a $\mathrm{U}(1)_Y \times \mathrm{SU}(2)_L
\times \mathrm{SU}(3)_C$ gauge field
$A_\mu$ from the manifold directions;
(ii)~a Higgs doublet $H \in (1,2,\tfrac{1}{2})$
from the finite directions.


%-----------------------------------------------------------------------
\subsection{Definition of $\Sigma_F$}
\label{sec:SigmaF}
%-----------------------------------------------------------------------

We define the internal QRE as
the Umegaki relative entropy between the
Gibbs state of the full (fluctuated) Dirac
operator and the Gibbs state of the free
(unfluctuated) Dirac operator:
\begin{equation}\label{eq:SigmaF_def}
  \boxed{\;
  \Sigma_F
  \;\equiv\;
  D\!\left(\rho_A \,\big\|\, \rho_0\right)
  \;}
\end{equation}
where
$\rho_A = e^{-\beta\, D_A^2}/ Z_A$
and
$\rho_0 = e^{-\beta\, D^{\,2}}/ Z_0$,
with $\beta = 1/\Lambda^2$.

By the standard identity for the QRE between
Gibbs states at the same $\beta$,
\begin{equation}\label{eq:SigmaF_expand}
  \Sigma_F
  = \beta\,\bigl\langle
    D^2 - D_A^2
  \bigr\rangle_{\!\rho_A}
  - \beta\,(F_A - F_0)\,,
\end{equation}
where $F = -\beta^{-1}\ln Z$ is the
free energy.
Substituting
$D_A^2 = D^2 + \{D,A_{\mathrm{in}}\}
+ A_{\mathrm{in}}^2$
(with $A_{\mathrm{in}} \equiv A + JAJ^{-1}$):
\begin{equation}\label{eq:SigmaF_inner}
  \Sigma_F
  = -\beta\,\bigl\langle
    \{D,A_{\mathrm{in}}\}
    + A_{\mathrm{in}}^2
  \bigr\rangle_{\!\rho_A}
  - \beta\,(F_A - F_0)\,.
\end{equation}


%-----------------------------------------------------------------------
\subsection{Variation of $\Sigma_F$ yields
the SM field equations}
\label{sec:variation}
%-----------------------------------------------------------------------

We show that
$\delta_A \Sigma_F = 0$ reproduces the
Euler--Lagrange equations of the
spectral action.
Using Eq.~\eqref{eq:SigmaF_expand} and
the CCSvS identification~\cite{CCSvS2018}
$\SvN(\rho_\beta)
= S_{\mathrm{spectral}}$,
the variation splits into
\begin{equation}\label{eq:var_split}
  \delta_A \Sigma_F
  = \beta\,\delta_A
    \bigl\langle D^2 \bigr\rangle_{\!\rho_A}
  - \delta_A S_{\mathrm{spectral}}\,.
\end{equation}
The first term is a boundary contribution
vanishing for compact~$M$.
Therefore
\begin{equation}\label{eq:equivalence}
  \delta_A \Sigma_F = 0
  \;\Longleftrightarrow\;
  \delta_A S_{\mathrm{spectral}} = 0\,,
\end{equation}
and the SM field equations follow
identically to the standard NCG
derivation~\cite{CC1996,vS2024}.


%-----------------------------------------------------------------------
\subsection{Fisher metric and the Weinberg angle}
\label{sec:Weinberg}
%-----------------------------------------------------------------------

The second-order expansion of $\Sigma_F$ in the
inner fluctuation defines the Fisher metric
on the gauge parameter space:
\begin{equation}\label{eq:FisherF}
  \Sigma_F(\varepsilon)
  = \frac{\varepsilon^2}{2}\,
    g_{AA}^{\mathrm{Fisher}}
  + \mathcal{O}(\varepsilon^3)\,.
\end{equation}
This metric decomposes along
$\mathfrak{u}(1) \oplus \mathfrak{su}(2)
\oplus \mathfrak{su}(3)$
with relative normalizations from the
quadratic Casimirs:
$g_3^2 = g_2^2
= \tfrac{5}{3}\, g_1^2$
at~$\Lambda$~\cite{CC1996,vS2024}.

The weak mixing angle acquires an
information-geometric interpretation:
\begin{equation}\label{eq:sin2theta}
  \sin^2\theta_W
  = \frac{g^{\mathrm{U}(1)}}
  {g^{\mathrm{U}(1)} + g^{\mathrm{SU}(2)}}
  = \frac{3}{8}
  \qquad\text{at } \Lambda\,,
\end{equation}
i.e., the fraction of
electroweak Fisher information carried by the
hypercharge sector.


%-----------------------------------------------------------------------
\subsection{Orthogonality of outer and inner sectors}
\label{sec:orthogonality}
%-----------------------------------------------------------------------

The total $\Sigma$ decomposes as
\begin{equation}\label{eq:Sigma_total}
  \Sigma_{\mathrm{total}}
  = \underbrace{2\ln Q}_{\text{outer, }a_2}
  \;+\;
  \underbrace{D(\rho_A \| \rho_0)}_{\text{inner, }a_4}\,.
\end{equation}
The cross-term of the Fisher metric
vanishes:
$g_{QA} = \partial^2 \Sigma_{\mathrm{total}}
/\partial Q\,\partial A |_{Q=1,A=0} = 0$,
because the $a_4$ sector is conformally
invariant in $d=4$~\cite{Wald1984}.

This means gravity and the gauge sector are
\emph{independent channels of entropy
production} at the Fisher level---a
decoupling specific to $d=4$.


%=======================================================================
\section{Discussion}
\label{sec:discussion}
%=======================================================================

\emph{Connection to Papers~1--6.}---Paper~1
established $\Sigma = -\ln F \geq 0$ as a
universal measure of temporal
asymmetry~\cite{PaperI}.
Paper~2 proved $\Sigma = 2\ln Q$ on static
backgrounds via the gravitational channel
theorem~\cite{PaperII}.
Paper~3 extended this to FRW backgrounds with
$Q = 1+\delta$ and resolved the
$c_s^2 = 0$ constraint~\cite{PaperIII}.
The present paper identifies the origin of the
``2'' in $\Sigma = 2\ln Q$: it is the
$(d-2)(d-3)$ Fisher information of the
Einstein--Hilbert spectral action, evaluated at
the physical spacetime dimension $d=4$.

\emph{Why $d=4$?}---Several arguments for
$d=4$ spacetime dimensions exist in the
literature: stability of planetary
orbits~\cite{Ehrenfest1917}, the Weyl
tensor's conformal properties, string theory
compactification, and the structure of division
algebras~\cite{BaezHuerta2010}.
The present argument is different: $d=4$ is
the unique dimension in which the conformal
Fisher information of the gravitational action
equals $2$, matching the entropy production
of the quantum channel associated with
temporal asymmetry~[NEW].
This is an algebraic condition---$(d-2)(d-3) = 2$---not
a dynamical selection principle.

\emph{The spectral action as entropy potential.}---The
CCSvS result~\cite{CCSvS2018} identifies the
von~Neumann entropy with the spectral action.
Our result identifies $\Sigma$ with the Fisher
information (second variation) of this entropy.
The spectral action is therefore the
``entropy potential'' whose curvature
in field space generates the temporal asymmetry
measure.
The Einstein--Hilbert term $a_2$ is the sector
whose curvature matches $\Sigma = 2\ln Q$;
the matter sector $a_4$ is flat (conformally
invariant) and does not contribute.

\emph{Towards the Standard Model.}---As shown in
Sec.~\ref{sec:AF}, the gauge and Higgs sectors
enter through inner fluctuations of the Dirac
operator on $\mathcal{A}_F$, defining
$\Sigma_F = D(\rho_A \| \rho_0)$ whose variation
reproduces the SM field equations.
The orthogonality
$g_{QA} = 0$ (Sec.~\ref{sec:orthogonality})
ensures that $\Sigma_{\mathrm{grav}}$ and
$\Sigma_F$ decouple at the Fisher level.
The weak mixing angle $\sin^2\theta_W = 3/8$
at~$\Lambda$ acquires a Fisher-information
interpretation (Sec.~\ref{sec:Weinberg}).
Further extensions to fermion masses require
the exceptional Jordan algebra
$J_3(\mathbb{O})$~\cite{Singh2025,FureyHughes2024}
and are the subject of Papers~10--11.

\emph{The Dorau--Much connection.}---Dorau and
Much~\cite{DorauMuch2025} recently derived the
semiclassical Einstein equations from QRE on
bifurcate Killing horizons.
Their result corresponds to the \emph{first-order}
variation of $\Sigma$ (the entropic force
$d\Sigma/dQ = 2/Q$), while our Fisher metric
is the \emph{second-order} variation.
The two are consistent and complementary:
Einstein's equations arise from $\delta\Sigma = 0$
at first order; the Fisher metric characterizes
the information geometry of the solution space
at second order.


%=======================================================================
\begin{acknowledgments}
The author thanks Mark~M.~Wilde for correspondence
on the Petz recovery map, and acknowledges the
foundational work of Alain Connes, Ali Chamseddine,
and Walter van Suijlekom on the spectral action
and its entropic interpretation.
\end{acknowledgments}
%=======================================================================


\begin{thebibliography}{99}

\bibitem{PaperI}
S.-K.~Huang,
``Temporal asymmetry from the Petz recovery map:
A unifying measure for the arrow of time,''
arXiv:2502.xxxxx (2026).

\bibitem{PaperII}
S.-K.~Huang,
``The gravitational refractive index as Petz
recovery fidelity,''
in preparation (2026).

\bibitem{PaperIII}
S.-K.~Huang,
``Dark matter as Khronon condensate: Weak-field
phenomenology of the $\tau$ framework,''
in preparation (2026).

\bibitem{Umegaki1962}
H.~Umegaki,
``Conditional expectation in an operator algebra.
IV. Entropy and information,''
K\=odai Math.\ Sem.\ Rep.\ \textbf{14}, 59 (1962).

\bibitem{CCSvS2018}
A.~H.~Chamseddine, A.~Connes, and
W.~D.~van~Suijlekom,
``Entropy and the spectral action,''
Commun.\ Math.\ Phys.\ \textbf{373}, 457 (2020);
arXiv:1809.02944.

\bibitem{DKvS2019}
R.~Dong, M.~Khalkhali, and
W.~D.~van~Suijlekom,
``Second quantization and the spectral action,''
J.\ Geom.\ Phys.\ \textbf{167}, 104285 (2021);
arXiv:1903.09624.

\bibitem{CC1996}
A.~H.~Chamseddine and A.~Connes,
``The spectral action principle,''
Commun.\ Math.\ Phys.\ \textbf{186}, 731 (1997);
hep-th/9606001.

\bibitem{CC2006}
A.~H.~Chamseddine and A.~Connes,
``Scale invariance in the spectral action,''
J.\ Math.\ Phys.\ \textbf{47}, 063504 (2006).

\bibitem{Vassilevich2003}
D.~V.~Vassilevich,
``Heat kernel expansion: User's manual,''
Phys.\ Rep.\ \textbf{388}, 279 (2003);
hep-th/0306138.

\bibitem{Gilkey1995}
P.~B.~Gilkey,
\emph{Invariance Theory, the Heat Equation, and
the Atiyah--Singer Index Theorem},
2nd ed.\ (CRC Press, Boca Raton, 1995).

\bibitem{Wald1984}
R.~M.~Wald,
\emph{General Relativity}
(University of Chicago Press, Chicago, 1984).

\bibitem{vS2024}
W.~D.~van~Suijlekom,
\emph{Noncommutative Geometry and Particle Physics},
2nd ed.\ (Springer, Cham, 2024).

\bibitem{BS2024}
L.~Blanchet and C.~Skordis,
``Dark matter as a Khronon condensate,''
arXiv:2404.06584 (2024).

\bibitem{DorauMuch2025}
A.~Dorau and A.~Much,
``From quantum relative entropy to the
semiclassical Einstein equations,''
Phys.\ Rev.\ Lett.\ (2025);
arXiv:2510.24491.

\bibitem{LVR2016}
N.~Lashkari and M.~Van~Raamsdonk,
``Canonical energy is quantum Fisher information,''
J.\ High Energy Phys.\ \textbf{04}, 153 (2016);
arXiv:1508.00897.

\bibitem{vandenDungen2016}
K.~van~den~Dungen,
``Krein spectral triples and the fermionic action,''
Math.\ Phys.\ Anal.\ Geom.\ \textbf{19}, 4 (2016).

\bibitem{Besnard2017}
F.~Besnard,
``On the definition of spacetimes in noncommutative
geometry: Part~I,''
J.\ Geom.\ Phys.\ \textbf{123}, 292 (2018);
arXiv:1611.07830.

\bibitem{Ehrenfest1917}
P.~Ehrenfest,
``In what way does it become manifest in the
fundamental laws of physics that space has three
dimensions?''
Proc.\ Amsterdam Acad.\ \textbf{20}, 200 (1917).

\bibitem{BaezHuerta2010}
J.~C.~Baez and J.~Huerta,
``Division algebras and supersymmetry~I,''
in \emph{Superstrings, Geometry, Topology,
and $C^*$-Algebras}, Proc.\ Symp.\ Pure Math.\
\textbf{81}, 65 (2010);
arXiv:0909.0551.

\bibitem{Singh2025}
T.~P.~Singh,
``Fermion mass ratios from the exceptional Jordan algebra,''
arXiv:2508.10131 (2025).

\bibitem{FureyHughes2024}
C.~Furey and N.~Hughes,
``Three generations and a trio of trialities,''
Phys.\ Lett.\ B (2025);
arXiv:2409.17948.

\bibitem{FarnsworthEtAl2026}
S.~Farnsworth, F.~Finster, N.~Paganini, and T.~P.~Singh,
``Causal fermion systems, non-commutative geometry
and generalized trace dynamics,''
arXiv:2603.05018 (2026).

\end{thebibliography}

\end{document}
